\label{subsec:analytic_geometry_equivalent_line_representations}

Eine Gerade besitzt keine eindeutige Parameterdarstellung. Das Erkennen
gleichwertiger Darstellungen verhindert, dass verschiedene Formeln
fälschlich als verschiedene Geraden interpretiert werden.

\begin{ExampleBox}{Dieselbe Gerade unterschiedlich darstellen}{analytic_geometry_equivalent_lines}
  Dieselbe Gerade lässt sich auf unterschiedliche Arten parametrisieren.

  \begin{center}
    \textcolor{olive!80!black}{$\overrightarrow{OP}$: Stützvektor}
    \qquad
    \textcolor{THMBlue}{$\vec a$: Richtungsvektor}
  \end{center}

  \noindent
  \begin{minipage}[t]{0.32\linewidth}
    \centering
    \begingroup
      \def\LineUnit{0.55}
      \def\LineSupportX{2}\def\LineSupportY{2}
      \def\LineDirectionX{1}\def\LineDirectionY{-1}
      \def\LineLambdaMin{-1.5}\def\LineLambdaMax{1.5}
      \def\LineShowX{0}
      \CachedTikzFigure{content/resources/chapters/analytic_geometry/tikz/parameterized_line_representation.tex}
    \endgroup
    \[
      g:\;\vec x=\begin{ClozeContinuation}[6]{2}
        \begin{pmatrix}2\\2\end{pmatrix}
        +\lambda\begin{pmatrix}1\\-1\end{pmatrix}
      \end{ClozeContinuation}
    \]
  \end{minipage}\hfill
  \begin{minipage}[t]{0.32\linewidth}
    \centering
    \begingroup
      \def\LineUnit{0.55}
      \def\LineSupportX{1}\def\LineSupportY{3}
      \def\LineDirectionX{1}\def\LineDirectionY{-1}
      \def\LineLambdaMin{-0.5}\def\LineLambdaMax{2.5}
      \def\LineShowX{0}
      \CachedTikzFigure{content/resources/chapters/analytic_geometry/tikz/parameterized_line_representation.tex}
    \endgroup
    \[
      g:\;\vec x=\begin{ClozeContinuation}[6]{2}
        \begin{pmatrix}1\\3\end{pmatrix}
        +\mu\begin{pmatrix}1\\-1\end{pmatrix}
      \end{ClozeContinuation}
    \]
  \end{minipage}\hfill
  \begin{minipage}[t]{0.32\linewidth}
    \centering
    \begingroup
      \def\LineUnit{0.55}
      \def\LineSupportX{2}\def\LineSupportY{2}
      \def\LineDirectionX{-1}\def\LineDirectionY{1}
      \def\LineLambdaMin{-1.5}\def\LineLambdaMax{1.5}
      \def\LineShowX{0}
      \CachedTikzFigure{content/resources/chapters/analytic_geometry/tikz/parameterized_line_representation.tex}
    \endgroup
    \[
      g:\;\vec x=\begin{ClozeContinuation}[6]{2}
        \begin{pmatrix}2\\2\end{pmatrix}
        +\nu\begin{pmatrix}-1\\1\end{pmatrix}
      \end{ClozeContinuation}
    \]
  \end{minipage}

  \par\bigskip

  Dieselbe Gerade entsteht, wenn
  \begin{itemize}[leftmargin=1.6em,itemsep=0.4em]
    \item ein anderer Punkt der Geraden als Aufpunkt gewählt wird und
    \item der Richtungsvektor durch ein beliebiges von null verschiedenes
      Vielfaches ersetzt wird.
  \end{itemize}
\end{ExampleBox}
